\documentclass[conference]{IEEEtran}
\usepackage{amsmath,url}
\newcommand{\toprule}{\hline}\newcommand{\midrule}{\hline}\newcommand{\bottomrule}{\hline}
\title{Version-Aware Evidence Grounding for Local Security Alert Explanations}
\author{\IEEEauthorblockN{Ananay Srivastava, Jitendra Singh, Manik Bhardwaj, Ashish Jain, Mukunj Singh}\IEEEauthorblockA{Department of Computer Science and Engineering, Sharda University, India}}
\begin{document}\maketitle
\begin{abstract}Local language models can make security alerts readable without exporting evidence, but fluent text can contain unsupported, stale, or unauthorized claims. We define a bounded local explanation pipeline that consumes a recipient-authorized evidence snapshot, separates observations, deductions, unknowns, and advice, and verifies claim references, versions, permissions, numeric bounds, and predicate relationships. Evidence correction, withdrawal, or authorization change invalidates dependent explanations; one repair is attempted before deterministic fallback. A finite eight-case study with a Qwen3-4B Q4_K_M CPU configuration measured templates, direct generation, cited generation, verification without lifecycle checks, and the complete method. Template output had full contract coverage; direct generation produced four unsupported claims across two cases, cited generation produced one unsupported claim and one generation failure, while the complete pipeline fell back in both cases after bounded verification. Median direct and cited generation times were 66.7 s and 95.0 s. Results establish a measurable safety pipeline and its current resource limit, not human usefulness or model superiority.\end{abstract}
\section{Introduction}The useful question for local AI in security operations is not whether a model can write fluent prose, but whether each factual claim is authorized and remains valid when evidence changes. We treat explanation as a typed admission problem over a versioned snapshot.
\section{Claim contract}An explanation contains four explicit sets: observations copied from evidence, deductions linked to predicates, unknowns, and advisory suggestions. Every factual claim references an evidence identifier and revision. The verifier rejects missing references, unauthorized identifiers, stale revisions, values outside snapshot bounds, and unsupported relationships. Advice is never presented as an observation. Model availability is independent of detection.
\section{Lifecycle protocol}The recipient receives an immutable snapshot hash and authorization scope. Generation is bounded to one concurrent request, fixed context/output limits, and non-thinking Q4\_K\_M inference. On correction, withdrawal, expiry, or permission change, dependent explanations become invalid and are removed from the usable set. A single repair prompt may run; failure produces a deterministic template containing only verified observations and unknowns.
\section{Evaluation}We froze eight claim cases containing benign alternatives, contradictory observations, stale revisions, unauthorized evidence, numeric bounds, and malicious evidence strings. We compared deterministic templates, direct local generation, citation-prompted generation, verification without lifecycle checks, and the complete method.\begin{table}[t]\centering\caption{Qwen3-4B CPU pilot (two cases per condition).}\begin{tabular}{lrrrr}\toprule Condition&Claims&Unsupported&Fallback&Median s\\\midrule Template&12&0&0&0.0003\\Direct&8&4&0&66.7\\Cited&4&1&0&95.0\\No lifecycle&12&0&2&--\\Complete&12&0&2&215.1\\\bottomrule\end{tabular}\end{table}The complete method rejected unsafe generations and used deterministic fallback in both cases. This is a safety result, not a quality score. No human evaluation was performed.
\section{Implementation and security}The local endpoint is interchangeable; the default configuration pins Qwen3-1.7B GGUF and evaluates Qwen3-4B GGUF. The API sends only authorized snapshots over an authenticated local connection. Prompt text is treated as hostile evidence, and generated identifiers are checked against the snapshot rather than trusted as citations. Dashboard records expose provenance, verification status, and service health without raw logs.
\section{Limitations}The pilot is small, CPU-bound, and measured on one machine. Two cases per model condition cannot support generalization. The complete pipeline's latency shows that bounded local inference may be unsuitable for interactive use without a faster device or smaller model. Grounding, claim verification, and revocation have adjacent prior art; the contribution is their explicit lifecycle contract and reproducible failure behavior in security explanations.
\section{Conclusion}Version-aware authorization turns local explanation into an auditable pipeline with safe fallback. The measured failures motivate larger frozen evaluations and human usefulness studies before deployment claims.
\begin{thebibliography}{00}\bibitem{hackphyr}Rigaki et al., ``Hackphyr: A local fine-tuned LLM agent for network security,'' arXiv:2409.11276, 2024.\bibitem{gsar}Dunietz et al., ``GSAR: Grounded security alert reasoning,'' arXiv:2604.23366, 2026.\bibitem{factscore}Min et al., ``FActScore,'' EMNLP, 2023.\bibitem{revocation}Liu et al., ``Revocation-aware retrieval for grounded generation,'' arXiv:2609.08258, 2026.\end{thebibliography}\end{document}
